\section*{Capítulo \ref{ch:b3:cytoskeleton-motility} --- Dinámica del citoesqueleto y motilidad celular}

\begin{solution}{exo:b3:cytoskeleton-motility:1}
Actina: \qty{7}{nm}, actina globular, polar (barbado/puntiagudo), ATP;
corteza, protrusión, contracción con la \omterm{def:b3:cytoskeleton-motility:motors}{miosina}. \omterm{def:b3:cytoskeleton-motility:filaments}{Microtúbulo}:
\qty{25}{nm}, dímero de $\alpha\beta$-tubulina, polar (más/menos), GTP;
transporte a larga distancia, huso, \omterm{def:b3:cytoskeleton-motility:cilia}{cilios}. \omterm{def:b3:cytoskeleton-motility:filaments}{Filamento intermedio}:
\qty{10}{nm}, proteínas en hélice enrollada (queratina, vimentina,
lamina), apolar, sin nucleótido; resistencia mecánica de la célula y del
núcleo.
\end{solution}

\begin{solution}{exo:b3:cytoskeleton-motility:2}
La concentración de monómero a la que un extremo ni crece ni se encoge,
$k_{\text{off}}/k_{\text{on}}$. Los dos extremos difieren porque la
subunidad hidroliza su ATP tras incorporarse, de modo que los extremos
presentan especies químicas distintas (actina-ATP que llega al extremo
barbado, actina-ADP que sale por el puntiagudo) con afinidades distintas.
Con concentraciones críticas iguales habría un solo equilibrio: ni noria,
ni flujo, ni manera de pagar un recambio dirigido.
\end{solution}

\begin{solution}{exo:b3:cytoskeleton-motility:3}
Periferia por \omterm{def:b3:cytoskeleton-motility:filaments}{microtúbulos}: \omterm{def:b3:cytoskeleton-motility:motors}{cinesina}, hacia el extremo más. Centro:
\omterm{def:b3:cytoskeleton-motility:motors}{dineína} citoplasmática, hacia el extremo menos. Fibra de tensión: \omterm{def:b3:cytoskeleton-motility:motors}{miosina}
II, hacia el extremo barbado de la actina. \omterm{def:b3:cytoskeleton-motility:cilia}{Cilio}: \omterm{def:b3:cytoskeleton-motility:cilia}{dineína axonémica},
hacia el extremo menos del doblete vecino (la base).
\end{solution}

\begin{solution}{exo:b3:cytoskeleton-motility:4}
Protrusión (Rac para los \omterm{def:b3:cytoskeleton-motility:migration}{lamelipodios}, Cdc42 para los filopodios),
adhesión (\omterm{def:b3:cytoskeleton-motility:migration}{integrinas}; por debajo de Rac y de Rho), contracción (Rho, a
través de la \omterm{def:b3:cytoskeleton-motility:motors}{miosina} II) y retracción de la parte trasera (la liberación
de las adhesiones, impulsada por la contracción).
\end{solution}

\begin{solution}{exo:b3:cytoskeleton-motility:5}
$C_{c}^{+} = 1/5 = \qty{0.2}{\micro M}$, $C_{c}^{-} = 0.8/1 =
\qty{0.8}{\micro M}$; $C_{\text{ss}} = (1 + 0.8)/(5 + 1) = \qty{0.3}{\micro
M}$; $J = 5\times 0.3 - 1 = 0.5$ subunidades por segundo,
\qty{1.35}{nm/s}.
\end{solution}

\begin{solution}{exo:b3:cytoskeleton-motility:6}
$800/8 = 100$ ATP por segundo. El metro cuesta $1/(8\times 10^{-7}) =
1.25\times 10^{6}$ s, unas dos semanas, y $1.25\times 10^{8}$ pasos y
ATP. La difusión tardaría \num{16000} años: mil millones de veces más.
\end{solution}

\begin{solution}{exo:b3:cytoskeleton-motility:7}
$F_{\max} = 8.3\times 10^{-20}/3.6\times 10^{-8} = \qty{2.3}{pN}$. La
misma energía repartida en un paso más largo da menos fuerza --- una
palanca más larga. La \omterm{def:b3:cytoskeleton-motility:motors}{cinesina} (paso corto, \qty{6}{pN}) conviene a las
cargas pesadas; la \omterm{def:b3:cytoskeleton-motility:motors}{miosina} V (paso largo) cubre más distancia por ATP y
conviene a un transporte rápido y ligero.
\end{solution}

\begin{solution}{exo:b3:cytoskeleton-motility:8}
(a) El taxol congela los \omterm{def:b3:cytoskeleton-motility:filaments}{microtúbulos}: el huso no puede buscar, capturar
ni corregir, y la mitosis se detiene en el punto de control; la célula
que repta sigue protruyendo pero pierde la polaridad a largo plazo.
(b) La colchicina elimina los \omterm{def:b3:cytoskeleton-motility:filaments}{microtúbulos}: sin huso, parada mitótica; el
reptar continúa sobre la actina pero sin dirección persistente. (c) La
citocalasina pone caperuza a los extremos barbados: sin protrusión, no
hay reptar; la mitosis avanza pero la citocinesis falla y quedan células
binucleadas. (d) Inhibición de Rac: sin \omterm{def:b3:cytoskeleton-motility:migration}{lamelipodios}, no hay reptar; la
división apenas se ve afectada.
\end{solution}

\begin{solution}{exo:b3:cytoskeleton-motility:9}
Tiempo de permanencia del orden del semitiempo, \qty{20}{s} (constante de
tiempo $20/\ln 2 \approx \qty{29}{s}$); fracción inmóvil, \qty{10}{\%}.
La polimerización en el borde ha de abastecer a la vez el avance y el
flujo retrógrado: $1 + 1.5 = \qty{2.5}{\micro m/min}$.
\end{solution}

\begin{solution}{exo:b3:cytoskeleton-motility:10}
Un extremo en crecimiento conserva una caperuza de tubulina-GTP; por
detrás de ella la hidrólisis produce tubulina-GDP tensada. Perder la
caperuza --- un suceso aleatorio cuya probabilidad sube cuando el
crecimiento se frena --- libera la tensión y el extremo se encoge deprisa
hasta que se forma una caperuza nueva. Justo por encima de la
\omterm{thm:b3:cytoskeleton-motility:treadmilling}{concentración crítica}, cada \omterm{def:b3:cytoskeleton-motility:filaments}{microtúbulo} está en cada momento o bien con
caperuza y creciendo o bien sin ella y desplomándose, de modo que las
longitudes divergen en una población larga y otra que se desvanece, en
lugar de agruparse en torno a una media. La célula gana una exploración
rápida de su volumen y la capacidad de estabilizar de forma selectiva ---
un extremo más que alcanza un cinetocoro o un punto cortical se captura y
se conserva, y los demás se reciclan en minutos: buscar y capturar.
\end{solution}

\begin{solution}{exo:b3:cytoskeleton-motility:11}
A esta escala la inercia es despreciable y las ecuaciones del fluido son
reversibles en el tiempo: un movimiento que se deshace a sí mismo (un
remo que va y vuelve) devuelve el cuerpo al punto de partida, sea cual
sea la velocidad de cada golpe. Una hélice que gira nunca se deshace a sí
misma --- su movimiento es quiral y continuo --- y por eso produce empuje
neto. A \qty{100}{Hz} con mil protones por vuelta, $10^{5}$ protones por
segundo y motor.
\end{solution}

\begin{solution}{exo:b3:cytoskeleton-motility:12}
Los \omterm{def:b3:cytoskeleton-motility:cilia}{cilios} de las vías respiratorias no pueden batir, el moco y las
bacterias no se aclaran y las infecciones se repiten. Los flagelos de los
espermatozoides, construidos sobre el mismo \omterm{def:b3:cytoskeleton-motility:cilia}{axonema}, son inmóviles:
infertilidad. En el embrión, los \omterm{def:b3:cytoskeleton-motility:cilia}{cilios} móviles del nódulo impulsan un
flujo hacia la izquierda que fija el eje izquierda--derecha; sin flujo el
lado se elige al azar, de modo que la mitad de los pacientes tienen los
órganos invertidos.
\end{solution}

\begin{solution}{pb:b3:cytoskeleton-motility:1}
\textbf{1.} $\qty{300}{\micro m^3} = 3\times 10^{-13}$ L; $\times
2\times 10^{-4}$ mol/L $= 6\times 10^{-17}$ mol $= 3.6\times 10^{7}$
moléculas; $1.8\times 10^{7}$ en filamentos.
\textbf{2.} $1.8\times 10^{7}\times \qty{2.7}{nm} = \qty{4.9}{cm}$ de
filamento; a \qty{0.25}{\micro m} cada uno, unos $2\times 10^{5}$
filamentos.
\textbf{3.} La proteína de caperuza bloquea la mayoría de los extremos
barbados, y la profilina y la timosina-$\beta$4 unen los monómeros de
modo que la actina-ATP verdaderamente libre está cerca de la
\omterm{thm:b3:cytoskeleton-motility:treadmilling}{concentración crítica}; el crecimiento se limita a los extremos que la
célula descubre.
\textbf{4.} $C_{\text{ss}} = 2.2/12.9 = \qty{0.17}{\micro M}$; $J = 11.6
\times 0.17 - 1.4 = 0.6$ subunidades por segundo; \qty{1.6}{nm/s} $=
\qty{0.1}{\micro m/min}$.
\textbf{5.} $1/0.1 = 10$ veces.
\textbf{6.} \qty{1}{\micro m/min} $= \qty{16.7}{nm/s} = 6.2$ subunidades
por segundo y filamento; $\times 2\times 10^{5} = 1.2\times 10^{6}$ ATP
por segundo, alrededor del $0.1\,\%$ del recambio de la célula.
\textbf{7.} $1/(8\times 10^{-7}) = 1.25\times 10^{6}$ s $\approx 14.5$
días; $1.25\times 10^{8}$ pasos y ATP.
\textbf{8.} $L^{2}/2D$: $(1)^{2}/(2\times 10^{-12}) = 5\times 10^{11}$ s
$\approx \num{16000}$ años para un metro; $(10^{-5})^{2}/(2\times
10^{-12}) = 50$ s para \qty{10}{\micro m}. La difusión sirve dentro de un
cuerpo celular; cualquier cosa más larga necesita motores.
\textbf{9.} $5\times 10^{4}/6.02\times 10^{23} = 8.3\times 10^{-20}$ J
$= 20\,k_{B}T$.
\textbf{10.} $8.3\times 10^{-20}/8\times 10^{-9} = \qty{10}{pN}$;
eficiencia $6/10 \approx 60\,\%$.
\textbf{11.} $F = 6\pi\times 0.01\times 10^{-7}\times 8\times 10^{-7} =
1.5\times 10^{-14}$ N $= \qty{0.015}{pN}$, cuatrocientas veces menos que
la \omterm{prop:b3:cytoskeleton-motility:physics}{fuerza de parada}: un solo motor basta y sobra contra la fricción
viscosa; las cargas reales son los obstáculos y los amarres.
\textbf{12.} $10^{4}\times 100 = 10^{6}$ ATP por segundo, el $0.1\,\%$
del presupuesto de la neurona: el transporte es barato. Pero los motores
se paran en cuanto cae el ATP, de modo que un fallo de las mitocondrias
priva a la sinapsis de todo lo que se fabrica en el cuerpo celular --- el
fallo del transporte axónico que se ve en las enfermedades
neurodegenerativas.
\textbf{13.} \qty{10}{\micro m/min} $= \qty{167}{nm/s}$; $167/2.7 = 62$
subunidades por segundo y filamento.
\textbf{14.} $200\times 20 = 4000$ filamentos; $4000\times 62 = 2.5\times
10^{5}$ subunidades por segundo.
\textbf{15.} $2.5\times 10^{5}$ ATP por segundo, el $0.025\,\%$ del
presupuesto.
\textbf{16.} $\qty{3}{\micro m}/(\qty{10}{\micro m/min}) = \qty{0.3}{min}
= \qty{18}{s}$; $2.5\times 10^{5}\times 18 = 4.5\times 10^{6}$
subunidades en la red.
\textbf{17.} $1\,\text{pN}\times 2.7\,\text{nm} = \qty{2.7}{pN.nm} =
0.66\,k_{B}T$: una fluctuación de esa energía ocurre con probabilidad
cercana a $e^{-0.66} \approx 0.5$ --- se abren huecos de una subunidad de
continuo, y el trinquete es del todo verosímil.
\textbf{18.} La propia polimerización empuja, siendo la bacteria la
``membrana'' del frente de su propia red; con la maquinaria Arp2/3 de la
célula alcanza la velocidad del \omterm{def:b3:cytoskeleton-motility:migration}{lamelipodio} y más --- hasta
\qty{0.5}{\micro m/s}, \qty{30}{\micro m/min}.
\textbf{19.} En $C = K_{d}$, $\theta = 1/2$ y $\mathrm{d}\theta/
\mathrm{d}C = 1/(4C)$: una diferencia del \qty{2}{\%} en $C$ da $\Delta\theta
= 0.005$; con \num{25000} receptores por lado, $125$ más ocupados en el
frente.
\textbf{20.} Cada lado tiene unos \num{12500} ocupados, que fluctúan en
$\sqrt{\num{12500}} \approx 110$, de modo que la diferencia entre los dos
lados fluctúa en unos $160$: los $125$ se pierden en el ruido en
cualquier instante. Los receptores se vuelven a unir alrededor de una vez por
segundo; promediar durante $N$ segundos encoge el ruido en $\sqrt{N}$ ---
diez segundos lo bajan a $50$ y el gradiente queda leído.
\textbf{21.} Rac (con Cdc42) en el frente, Rho en la parte trasera. Con
Rac activa por todas partes la célula protruye en todas direcciones, se
extiende y no va a ninguna parte.
\textbf{22.} $30/3 = \qty{10}{\micro m/min}$: tal como se daba.
\textbf{23.} La protrusión se sustituye por ampollas impulsadas por la
presión; la adhesión y la contracción siguen dependiendo de la actina y
la \omterm{def:b3:cytoskeleton-motility:motors}{miosina}, de modo que la actina sigue siendo esencial para el resto del
ciclo.
\textbf{24.} El frente se reconstruye cada \qty{18}{s}: un cambio en la
actividad de Rac redirige la polimerización en una sola vida de la red,
el nuevo frente encabeza la marcha y los pasos lentos de la parte trasera
lo siguen.
\textbf{25.} Noria in vitro, \qty{1.6}{nm/s}; unos $14.5$ días para
recorrer el axón; y unos $2.5\times 10^{5}$ ATP por segundo para la
protrusión.
\end{solution}
